\subsection{Classes of equivalent functions}

Let \(\mu\) be a measure on a locally compact space \(X\) . Given a set \(F\) , two mappings \(f, g\) of \(X\) into \(F\) are said to be equivalent with respect to \(\mu\) (or \(\mu\) -equivalent, or simply equivalent if no confusion can arise) if \(f(x) = g(x)\) almost everywhere in \(X\) . Since the union of two negligible sets is negligible, one indeed defines in this way an equivalence relation in the set \(F^X\) of all mappings of \(X\) into \(F\) ; when we speak of the equivalence class of such a function \(f\) (without further specification) it will be understood to be the class of all the functions equal almost everywhere to \(f\) ; in this chapter and those that follow, we will indicate this class by the notation \(\widetilde{f}\) .

PROPOSITION 8. --- Let \((\mathrm{F}_{n})\) be a countable family (finite or infinite) of sets. For every index n, let \(f_{n}, g_{n}\) be two equivalent mappings of X into \(F_{n}\) ; then, there exists a negligible set H such that, for every \(x \notin H\) , \(f_{n}(x) = g_{n}(x)\) for all n.

For, the set \(H_{n}\) of \(x \in X\) such that \(f_{n}(x) \neq g_{n}(x)\) is negligible, therefore so is their union H (No.~2, Prop. 4), and this set meets the requirements.

COROLLARY. --- If \(\varphi\) is a mapping of \(\prod_{n} F_n\) into a set \(G\) , then the mappings \(\varphi((f_n))\) and \(\varphi((g_n))\) of \(X\) into \(G\) are equivalent.

We denote by \(\varphi((\widetilde{f}_n))\) the equivalence class of every function \(\varphi((f_n))\) , where \(f_n\) is an arbitrary function in the class \(\widetilde{f}_n\) .

In particular, if F is a vector space over R, one defines \(\widetilde{f} + \widetilde{g}\) and \(\alpha\widetilde{f}\) to be the equivalence classes of \(f + g\) and \(\alpha f\) , respectively (f and g being mappings of X into F, and \(\alpha\) a scalar); we obtain in this way, on the set of equivalence classes of mappings of X into F, a vector space structure: moreover, this is the quotient space structure of that of \(F^{X}\) by the linear subspace of mappings f such that \(\widetilde{f} = \widetilde{0}\) (the functions that are zero almost everywhere), which we also call negligible functions (with values in F). One defines similarly the product \(\widetilde{g}\widetilde{f}\) , where \(\widetilde{f}\) is an equivalence class of mappings of X into F, and \(\widetilde{g}\) is an equivalence class of (finite) numerical functions defined on X: the set of equivalence classes of mappings of X into F is thus equipped with the structure of a module over the set of equivalence classes of finite numerical functions defined on X (which is itself equipped with a ring structure). If F is an algebra over R, one defines similarly an algebra structure on the set of equivalence classes of mappings of X into F.

Let F be a metrizable topological space, and consider a uniform structure on F compatible with its topology and defined by a countable family of pseudometrics \(\rho_{n}\) (GT, IX, §§1 and 2); in order that two mappings f, g of X into F be equivalent, it is necessary and sufficient that the numerical functions \(\rho_{n}(f,g)\) be negligible; for, this is equivalent to saying that there exists a negligible set H in X such that, for every \(x\notin H\) , \(\rho_{n}(f(x),g(x))=0\) for all n, that is, \(f(x)=g(x)\) . In particular, if F is a metrizable locally convex space, and \((q_{n})\) is a countable family of seminorms defining the topology of F (TVS, II, §4, No.~1), in order that two mappings f, g of X into F be equivalent it is necessary and sufficient that all of the numerical functions \(q_{n}(\mathbf{f}(x)-\mathbf{g}(x))\) be negligible.

PROPOSITION 9. --- Let f and g be two continuous mappings of X into a Hausdorff topological space F; for f and g to be equivalent, it is necessary and sufficient that \(f(x) = g(x)\) at every point of the support of \(\mu\) .

For, the set of \(x \in X\) such that \(f(x) \neq g(x)\) is an open set (GT, I, §8, No.~1); for it to be negligible, it is necessary and sufficient that it not intersect the support of \(\mu\) (No.~2, Prop. 5).

PROPOSITION 10. --- Let F be a Hausdorff locally convex space over R such that there exists in the dual F' of F a sequence ( \(a_{n}^{\prime}\) ) that is dense for the weak topology \(\sigma(F^{\prime},F)\) (TVS, II, §6, No.~2). In order that two mappings f, g of X into F be equivalent, it is necessary and sufficient that, for every n, the numerical functions \(\langle\mathbf{f}(x),\mathbf{a}_{n}^{\prime}\rangle\) and \(\langle\mathbf{g}(x),\mathbf{a}_{n}^{\prime}\rangle\) be equivalent.

The condition is obviously necessary. Conversely, if it is satisfied, there exists a negligible set H such that, for each \(x \notin H\) , \(\langle \mathbf{f}(x), \mathbf{a}_{n}^{\prime} \rangle = \langle \mathbf{g}(x), \mathbf{a}_{n}^{\prime} \rangle\) for every n; this means that the weakly continuous linear forms \(\mathbf{z}^{\prime} \mapsto \langle \mathbf{f}(x), \mathbf{z}^{\prime} \rangle\) and \(\mathbf{z}^{\prime} \mapsto \langle \mathbf{g}(x), \mathbf{z}^{\prime} \rangle\) on \(F^{\prime}\) are equal at each of the points \(a_{n}^{\prime}\) , hence are identical by virtue of the hypothesis, which proves that \(\mathbf{f}(x) = \mathbf{g}(x)\) for all \(x \notin H\) .

Note that the hypothesis of Prop. 10 is applicable in particular when F is a locally convex space that is metrizable and separable \(^{1}\) (TVS, III, §3, No.~4, Cor. 2 of Prop. 6).

